%%%%%%%%%%%%%%%%
%%% lev 10 %%%
%%%%%%%%%%%%%%%%

\Chapter{10}

\Verse{1}
{
	\hP{וַיִּקְחוּ בְנֵי אַהֲרֹן נָדָב וַאֲבִיהוּא אִישׁ מַחְתָּתוֹ וַיִּתְּנוּ בָהֵן אֵשׁ וַיָּשִׂימוּ עָלֶיהָ קְטֹרֶת וַיַּקְרִיבוּ לִפְנֵי יְהֹוָה אֵשׁ זָרָה אֲשֶׁר לֹא צִוָּה אֹתָם׃}
}
{
	\eP{
		And Aaron’s sons, Nadab and Abihu, each took his
		fire-holder, and they put fire in them and set incense on
		it. And they brought forward unfitting fire, which He had
		not commanded them, in front of YHWH.
	}
}
{
%	\footnote{
%		Nadab and Abihu. These are parallel to the names Nadab and Abiyah, the
%		sons of Jeroboam I, the first king of the northern kingdom of Israel (1
%		Kings 14:1; 15:25). King Jeroboam erects golden calves at sanctuaries at
%		Dan and Beth-El. The similarity of the names of the sons of the two men
%		who make golden calves is striking (Exodus 32; 1 Kings 12:28). Footnote
%		10:1. unfitting fire, which He had not commanded them. It is not clear
%		just what Nadab’s and Abihu’s offense is. The text says that they bring
%		fire that is zrh, that YHWH had not commanded them. Some have taken this
%		to mean “foreign” fire and therefore have related Nadab’s and Abihu’s
%		act to foreign worship, i.e., to another god. The term zrh, however,
%		never has the meaning of “foreign” in any other occurrence in the
%		Torah’s law or narrative. It rather denotes a person or thing that is
%		outside a particular group. In any context dealing with priests, for
%		example, it refers to a layperson. In the only other context referring
%		to incense burning (Exod 30:9), “unfitting” incense is listed along with
%		sacrifices and libations as things that should not be brought on the
%		special incense altar located inside the Tent of Meeting. That is, there
%		is no suggestion that it means anything foreign, but simply something
%		that is outside the realm of what is permitted, or, as our text says
%		explicitly, “that YHWH had not commanded.” Nadab and Abihu might thus be
%		thought of as having had positive intentions—to bring an incense
%		offering to God—but as unfortunately having acted incorrectly. Or,
%		alternatively, we might imagine a variety of less commendable motives in
%		them. The point is that the text does not deal with their motives
%		because that is not the issue. In the realm of the ritual, they have
%		failed to observe a boundary, and so their fate is settled. This is one
%		of several biblical stories that indicate that on the highest levels of
%		the ritual realm, intention does not matter. In the ethical realm it
%		does. The killing of a man, for example, may be found to be murder or to
%		be manslaughter, depending in part on the intention of the person who
%		killed him. In the ritual category, however, there are cases in which
%		innocent motives still do not make one innocent. Only a priest, for
%		example, can enter the Tent of Meeting. If a layperson enters, that
%		person must die, and the law does not provide for taking the motives for
%		the trespass into account. And in the famous case of the ark in 2 Samuel
%		6, Uzzah touches the ark to steady it because it is rocking on the backs
%		of oxen. His motive is good. But he dies because, like Nadab and Abihu,
%		he has violated the boundary of the holy.
%	}
}

\Verse{2}
{
	\hP{וַתֵּצֵא אֵשׁ מִלִּפְנֵי יְהֹוָה וַתֹּאכַל אוֹתָם וַיָּמֻתוּ לִפְנֵי יְהֹוָה׃}
}
{
	\eP{
		And fire came out from in front of YHWH and consumed them!
		And they died in front of YHWH.
	}
}
{
%	\footnote{
%		fire came out from in front of YHWH and consumed. Aaron’s inaugural day
%		of priestly sacrifice is thus to have been a glorious day of pomp and
%		ceremony, miraculously sanctioned by the divine glory and fire; but then
%		this extraordinary thing happens. The two eldest of his four sons bring
%		this offering that YHWH had not told them to make, and the divine fire
%		comes and consumes them. After seven chapters of laws of sacrifices and
%		two chapters describing the ordination rituals, the sudden account of
%		this horrible event in the middle of the ceremonies comes as a shock. It
%		is even more powerful because the exact same words that describe the
%		miraculous consumption of the ordination sacrifice two verses earlier
%		(9:24) now describe the miraculous killing of Nadab and Abihu. It is
%		possible that this account of the deaths of Nadab and Abihu should be
%		regarded as a separate story from the account of the consecration. It is
%		frequently understood that way, especially since it occurs here at the
%		beginning of a new chapter. Still, the division of the biblical books
%		into chapters was made centuries after composition of the books, and
%		there is no break in the flow of the story from the consecration of the
%		priesthood to the death of Aaron’s sons in the manuscripts of Leviticus.
%		Moreover, the last verses of the story of the consecration picture the
%		events as happening in the presence of the entire people, who are
%		assembled (9:22–24); and Moses explicitly refers to the presence of the
%		people in his first words to Aaron after his sons’ deaths. He says,
%		“That is what YHWH spoke, saying, ‘I shall be made holy through those
%		who are close to me, and I shall be honored in front of all the people’”
%		(10:3). The story of Nadab and Abihu is therefore most probably to be
%		understood as a horror that occurs during the inauguration of the
%		priesthood.
%	}
}

\Verse{3}
{
	\hP{וַיֹּאמֶר מֹשֶׁה אֶל אַהֲרֹן הוּא אֲשֶׁר דִּבֶּר יְהֹוָה לֵאמֹר בִּקְרֹבַי אֶקָּדֵשׁ וְעַל פְּנֵי כל הָעָם אֶכָּבֵד וַיִּדֹּם אַהֲרֹן׃}
}
{
	\eP{
		And Moses said to Aaron, “That is what YHWH spoke, saying,
		‘I shall be made holy through those who are close to me,
		and I shall be honored in front of all the people.’” And
		Aaron was silent.
	}
}
{
%	\footnote{
%		That is what YHWH spoke. The meaning of Moses’ remark to Aaron is
%		another provocative case of ambiguity in the Torah’s narrative. Is Moses
%		to be viewed as saying these words in horrified recognition of a truth?
%		or as a stinging reprimand to Aaron concerning his sons’ behavior? or as
%		an explanation? The text here specifically notes that Aaron makes no
%		response (“And Aaron was silent”), which is understandable but not
%		helpful and is itself open to a gamut of interpretations. Footnote 10:3.
%		I shall be made holy through those who are close to me. The text goes on
%		conveying that there is a burden on “those who are close to me,” as it
%		describes the removal of Aaron’s sons’ bodies by their cousins, and as
%		Moses instructs Aaron and his remaining sons, Eleazar and Ithamar, that
%		they must not mourn for Nadab and Abihu. There follows an account that
%		further conveys the pain within the priestly family. Moses finds that
%		Aaron and his remaining sons have not eaten the sacrifices of the day as
%		commanded, and he reprimands Eleazar and Ithamar angrily. Aaron replies
%		that, after all that has happened on this day, “If I had eaten a sin
%		offering today would it be good in YHWH’s eyes?” And the story concludes
%		with a play on this last wording, noting that Aaron’s answer “was good
%		in Moses’ eyes” (10:19–20). Their pain is a reminder that the standard
%		for leaders is tougher than for others. According to the Torah, leaders
%		do not get away with more because of their positions. Priests, prophets,
%		kings, rabbis, presidents: they suffer harder consequences. This in turn
%		sets up what will happen to Moses and to Aaron themselves later. Moses
%		quotes God here as saying, “I shall be made holy through those who are
%		close to me.” For those of us who know what is coming later in the
%		story, this is a shivering preparation for the episode of Moses’ own
%		sin. When Moses will strike the rock at Meribah, God will impose a
%		frightful consequence for him and for Aaron. And the reason that God
%		will give to Moses and Aaron will be: “Because you did not trust in me,
%		to make me holy before the eyes of the children of Israel” (Num 20:12).
%		Moses’ own words to Aaron here will come to testify against him there.
%		And as Aaron is silent here, Moses makes no answer there.
%	}
}

\Verse{4}
{
	\hP{וַיִּקְרָא מֹשֶׁה אֶל מִישָׁאֵל וְאֶל אֶלְצָפָן בְּנֵי עֻזִּיאֵל דֹּד אַהֲרֹן וַיֹּאמֶר אֲלֵהֶם קִרְבוּ שְׂאוּ אֶת אֲחֵיכֶם מֵאֵת פְּנֵי הַקֹּדֶשׁ אֶל מִחוּץ לַמַּחֲנֶה׃}
}
{
	\eP{
		And Moses called to Mishael and to Elzaphan, sons of
		Uzziel, Aaron’s uncle, and said to them, “Come forward.
		Carry your brothers from in front of the Holy to the
		outside of the camp.”
	}
}
{
}

\Verse{5}
{
	\hP{וַיִּקְרְבוּ וַיִּשָּׂאֻם בְּכֻתֳּנֹתָם אֶל מִחוּץ לַמַּחֲנֶה כַּאֲשֶׁר דִּבֶּר מֹשֶׁה׃}
}
{
	\eP{
		And they came forward and carried them by their coats to
		the outside of the camp as Moses had spoken.
	}
}
{
}

\Verse{6}
{
	\hP{וַיֹּאמֶר מֹשֶׁה אֶל אַהֲרֹן וּלְאֶלְעָזָר וּלְאִיתָמָר בָּנָיו רָאשֵׁיכֶם אַל תִּפְרָעוּ וּבִגְדֵיכֶם לֹא תִפְרֹמוּ וְלֹא תָמֻתוּ וְעַל כּל הָעֵדָה יִקְצֹף וַאֲחֵיכֶם כּל בֵּית יִשְׂרָאֵל יִבְכּוּ אֶת הַשְּׂרֵפָה אֲשֶׁר שָׂרַף יְהֹוָה׃}
}
{
	\eP{
		And Moses said to Aaron and to Eleazar and to Ithamar, his
		sons, “Don’t let loose the hair of your heads and don’t
		tear your clothes, so you won’t die and He’ll be angry at
		all the congregation. And your brothers, all the house of
		Israel, will weep for the burning that YHWH has made.
	}
}
{
}

\Verse{7}
{
	\hP{וּמִפֶּתַח אֹהֶל מוֹעֵד לֹא תֵצְאוּ פֶּן תָּמֻתוּ כִּי שֶׁמֶן מִשְׁחַת יְהֹוָה עֲלֵיכֶם וַיַּעֲשׂוּ כִּדְבַר מֹשֶׁה׃}
}
{
	\eP{
		And you shall not go out from the entrance of the Tent of
		Meeting, or else you’ll die, because YHWH’s anointing oil
		is on you.” And they did according to Moses’ word.
	}
}
{
}

\Verse{8}
{
	\hP{וַיְדַבֵּר יְהֹוָה אֶל אַהֲרֹן לֵאמֹר׃}
}
{
	\eP{And YHWH spoke to Aaron, saying,}
}
{
%	\footnote{
%		YHWH spoke to Aaron. This is the first time that God speaks directly to
%		Aaron alone since He first sent him to meet Moses (Exod 4:27), and it is
%		right after the death of his sons! It may be that the significance of
%		what Nadab and Abihu have done is the reason that the deity now
%		addresses Aaron directly concerning the limitations and responsibilities
%		of his family. Or perhaps we should understand this as an act of comfort
%		from God to Aaron after his frightful loss.
%	}
}

\Verse{9}
{
	\hP{יַיִן וְשֵׁכָר אַל תֵּשְׁתְּ אַתָּה וּבָנֶיךָ אִתָּךְ בְּבֹאֲכֶם אֶל אֹהֶל מוֹעֵד וְלֹא תָמֻתוּ חֻקַּת עוֹלָם לְדֹרֹתֵיכֶם׃}
}
{
	\eP{
		“You shall not drink wine and beer, you and your sons with
		you, when you come to the Tent of Meeting, so you won’t
		die. It is an eternal law through your generations,
	}
}
{
%	\footnote{
%		You shall not drink wine and beer. Some have suggested on the basis of
%		this verse that Nadab and Abihu had been drunk and that this is what
%		caused them to commit the offense. Admittedly, this prohibition of
%		alcohol for the priests when they enter the Tabernacle seems otherwise
%		unrelated to the context, but that is not sufficient grounds to
%		denigrate these priests. Whether one believes them to be historical
%		persons or literary characters, one must be circumspect in judging
%		biblical figures. It is good practice for being circumspect when we
%		judge living persons.
%	}
}

\Verse{10}
{
	\hP{וּלְהַבְדִּיל בֵּין הַקֹּדֶשׁ וּבֵין הַחֹל וּבֵין הַטָּמֵא וּבֵין הַטָּהוֹר׃}
}
{
	\eP{
		and to distinguish between the holy and the secular, and
		between the impure and the pure,
	}
}
{
%	\footnote{
%		to distinguish. Leviticus is concerned with orderliness. This
%		orderliness is reminiscent of the creation account in Genesis 1. There
%		are key parallels of wording, especially the term for distinction
%		(lhabdîl). As God creates by making distinctions, expressed in divine
%		speech, in Genesis (“God distinguished between the light and the
%		darkness”), so the function of the priesthood is described in Leviticus
%		as “to distinguish (lhabdîl) between the holy and the secular, and
%		between the impure and the pure.” Here law is conceived as a reflection
%		in the human realm of the order that was originally pictured in the
%		cosmic realm. Footnote 10:10. to distinguish between the holy and the
%		secular. Rabbi Simcha Weiser quoted an expression to me: The problem
%		with American Jews today is that they know how to make kiddush (the
%		blessing over wine), but they don’t know how to make habdalah (the
%		ceremony ending the Sabbath and beginning the secular week). That is,
%		they do not know how to distinguish between the holy and the secular.
%		James Kugel sent me something he wrote: “One of the crucial concepts of
%		biblical religion is almost altogether absent from modern life. The
%		whole of the biblical world was divided into two great domains, the holy
%		and the profane.” In The Hidden Face of God, I hinted at a path through
%		discoveries of the last century about the origin of the universe that
%		may lead us back to appreciation of the holy, the awesome, the wonder of
%		the universe. Footnote 10:10. between the impure and the pure. The
%		Hebrew terms here are also frequently translated as the “unclean” and
%		the “clean.” Neither pair of terms quite conveys the sense of the
%		Hebrew, which, first of all, employs terms, m and hôr, having two
%		different roots, thus perhaps conveying more clearly that these terms
%		refer to two distinct conditions, each of which has a particular legal
%		status. That is, being m is more than merely being un-hôr. To be m is to
%		be in a particular condition, the qualities of which are invisible but
%		which can be transmitted to other persons or objects by contact.
%		Therefore, physical separation and various prescribed acts are required
%		to remove this condition. One comes to be in a condition of being m as a
%		result of specific occurrences: from menstruation, from other flows of
%		matter from the body’s sexual organs, from a woman’s giving birth (for a
%		week if it is a male baby, two weeks if it is a female), from contact
%		with a corpse (of a human, of a forbidden animal, or of a permitted
%		animal if it died of disease or was killed by another animal), from
%		certain forms of leprosy, or from contact with persons or objects that
%		are m. (One may also become m from a forbidden sexual relationship,
%		although the term may possibly be understood differently in the relevant
%		passage, 18:24–30, and is in fact usually rendered differently in
%		translations; cf. also the case of human sacrifice, Footnote 20:3.) When
%		one performs certain acts to restrict and remove this condition
%		(separation, bathing, and, in some circumstances, sacrifice), one
%		returns to a condition of being hôr. The common element appears to be
%		that all these cases that make an individual m’ involve some sort of
%		visible change in body fluids and/or skin, although here, as with the It
%		is sometimes suggested that all these cases are related to death. But
%		this is not correct. One must account for the inclusion of menstruation
%		and childbirth, which relate to the start of life, not to its end. We
%		may extend the explanation, then, and say that all these cases relate to
%		the start or the end of life. But even this does not really account for
%		the inclusion of leprosy among the m conditions. My colleague Jacob
%		Milgrom has suggested that the leprous conditions all involve the
%		wasting of a body in some way that resembles the wasting of a corpse.
%		This is possible but hard to prove in the absence of evidence that
%		people in biblical times made this connection in their minds. The
%		evidence that Milgrom gives is the case in which Miriam is stricken with
%		leprosy, and Aaron pleads with Moses, “Let her not be like the dead who,
%		when he comes out of his mother’s womb, half of his flesh is eaten up!”
%		But this is a unique case: Miriam is stricken by God with a kind of
%		leprosy so that she is rendered utterly pale (“leprous as snow”). This
%		singular colorless condition of her skin may be what makes her
%		comparable to a corpse. And Aaron is comparing her not just to any
%		corpse but to a stillbirth. His exclamation in that story is not
%		adequate evidence to make a general connection between leprosy and
%		death. The essential common point, it seems to me, is that these are all
%		things that most people in most societies instinctively do not want to
%		touch: blood, semen, diseased skin, corpses. It may be that ancient
%		Israelites came to give all these things that they found repugnant to
%		touch a common name: m’. And then, in a second stage, people came to
%		think of m’ as a category and felt the need to find a common underlying
%		factor. Was being m’ understood to be a condition in which one was
%		covered with some sort of microscopic creatures like a spreading
%		bacterium that had to be washed away (as some have believed to this
%		day), or was it conceived as comparable to an infectious illness, or was
%		it strictly a status that was conceived in legal definition? Biblical
%		law, like biblical narrative, has its share of ambiguities, albeit of a
%		different sort from narrative ambiguity. These gaps in the information
%		that the text provides exist, first, because the laws are pictured as
%		coming directly from God, who is not bound to explain or justify His
%		commandments to those whom He commands. Second, the law codes, like
%		constitutional law, are statements of the primary principles and cases,
%		leaving the details and variations to be worked out by priests or judges
%		as they arise (Lev 10:10–11). Third, much of the law was practiced in
%		the daily life of the people of Israel over centuries, and so what
%		appear as gaps to us may have been common knowledge in the community in
%		which the text was composed.
%	}
}

\Verse{11}
{
	\hP{וּלְהוֹרֹת אֶת בְּנֵי יִשְׂרָאֵל אֵת כּל הַחֻקִּים אֲשֶׁר דִּבֶּר יְהֹוָה אֲלֵיהֶם בְּיַד מֹשֶׁה׃}
}
{
	\eP{
		and to instruct the children of Israel all the laws that
		YHWH has spoken to them by the hand of Moses.”
	}
}
{
}

\Verse{12}
{
	\hP{וַיְדַבֵּר מֹשֶׁה אֶל אַהֲרֹן וְאֶל אֶלְעָזָר וְאֶל אִיתָמָר בָּנָיו הַנּוֹתָרִים קְחוּ אֶת הַמִּנְחָה הַנּוֹתֶרֶת מֵאִשֵּׁי יְהֹוָה וְאִכְלוּהָ מַצּוֹת אֵצֶל הַמִּזְבֵּחַ כִּי קֹדֶשׁ קדָשִׁים הִוא׃}
}
{
	\eP{
		And Moses said to Aaron and to Eleazar and to Ithamar, his
		sons who were left, “Take the grain offering that is left
		from YHWH’s offerings by fire and eat it as unleavened
		bread beside the altar, because it is the holy of holies.
	}
}
{
}

\Verse{13}
{
	\hP{וַאֲכַלְתֶּם אֹתָהּ בְּמָקוֹם קָדוֹשׁ כִּי חקְךָ וְחק בָּנֶיךָ הִוא מֵאִשֵּׁי יְהֹוָה כִּי כֵן צֻוֵּיתִי׃}
}
{
	\eP{
		And you shall eat it in a holy place, because it is your
		statutory share, and it is your sons’ statutory share from
		YHWH’s offerings by fire, because that is what I was
		commanded.
	}
}
{
}

\Verse{14}
{
	\hP{וְאֵת חֲזֵה הַתְּנוּפָה וְאֵת שׁוֹק הַתְּרוּמָה תֹּאכְלוּ בְּמָקוֹם טָהוֹר אַתָּה וּבָנֶיךָ וּבְנֹתֶיךָ אִתָּךְ כִּי חקְךָ וְחק בָּנֶיךָ נִתְּנוּ מִזִּבְחֵי שַׁלְמֵי בְּנֵי יִשְׂרָאֵל׃}
}
{
	\eP{
		And you shall eat the breast of the elevation offering and
		the thigh of the donation in a pure place—you and your
		sons and your daughters with you—because they have been
		given as your statutory share and your sons’ statutory
		share from the peace-offering sacrifices of the children
		of Israel.
	}
}
{
}

\Verse{15}
{
	\hP{שׁוֹק הַתְּרוּמָה וַחֲזֵה הַתְּנוּפָה עַל אִשֵּׁי הַחֲלָבִים יָבִיאוּ לְהָנִיף תְּנוּפָה לִפְנֵי יְהֹוָה וְהָיָה לְךָ וּלְבָנֶיךָ אִתְּךָ לְחק עוֹלָם כַּאֲשֶׁר צִוָּה יְהֹוָה׃}
}
{
	\eP{
		They shall bring the thigh of the donation and the breast
		of the elevation offering with the offerings by fire of
		the fats to elevate, an elevation offering in front of
		YHWH; and it will be yours and your sons’ with you as an
		eternal law, as YHWH commanded.”
	}
}
{
}

\Verse{16}
{
	\hP{וְאֵת שְׂעִיר הַחַטָּאת דָּרֹשׁ דָּרַשׁ מֹשֶׁה וְהִנֵּה שֹׂרָף וַיִּקְצֹף עַל אֶלְעָזָר וְעַל אִיתָמָר בְּנֵי אַהֲרֹן הַנּוֹתָרִם לֵאמֹר׃}
}
{
	\eP{
		And Moses inquired about the goat of the sin offering, and
		here it was burnt! And he was angry at Eleazar and at
		Ithamar, Aaron’s sons who were left, saying,
	}
}
{
%	\footnote{
%		inquired. This is the middle of the Torah in terms of words, and,
%		appropriately, it is the word “to inquire” (Hebrew ), the word that
%		later comes to mean exegesis of the Torah. (See the comment on Footnote
%		11:42.)
%	}
}

\Verse{17}
{
	\hP{מַדּוּעַ לֹא אֲכַלְתֶּם אֶת הַחַטָּאת בִּמְקוֹם הַקֹּדֶשׁ כִּי קֹדֶשׁ קדָשִׁים הִוא וְאֹתָהּ נָתַן לָכֶם לָשֵׂאת אֶת עֲוֺן הָעֵדָה לְכַפֵּר עֲלֵיהֶם לִפְנֵי יְהֹוָה׃}
}
{
	\eP{
		“Why didn’t you eat the sin offering in the place of the
		Holy—because it’s the holy of holies, and He gave it to
		you for bearing the congregation’s crime, for making
		atonement over them in front of YHWH!
	}
}
{
}

\Verse{18}
{
	\hP{הֵן לֹא הוּבָא אֶת דָּמָהּ אֶל הַקֹּדֶשׁ פְּנִימָה אָכוֹל תֹּאכְלוּ אֹתָהּ בַּקֹּדֶשׁ כַּאֲשֶׁר צִוֵּיתִי׃}
}
{
	\eP{
		Here, its blood hasn’t been brought to the Holy, inside.
		You should have eaten it in the Holy as I commanded!”
	}
}
{
%	\footnote{
%		as I commanded. See the comment on Lev 8:31. Here again Moses speaks of
%		himself as commanding something concerning Aaron. And here again there
%		is uncertainty about the text, since other versions differ from this
%		reading in the Masoretic Text; they read that it is God who commands
%		Moses. The uncertainty in the technical matter of the differences in the
%		manuscripts underlines the larger issue: the growth in the stature and
%		authority of Moses.
%	}
}

\Verse{19}
{
	\hP{וַיְדַבֵּר אַהֲרֹן אֶל מֹשֶׁה הֵן הַיּוֹם הִקְרִיבוּ אֶת חַטָּאתָם וְאֶת עֹלָתָם לִפְנֵי יְהֹוָה וַתִּקְרֶאנָה אֹתִי כָּאֵלֶּה וְאָכַלְתִּי חַטָּאת הַיּוֹם הַיִּיטַב בְּעֵינֵי יְהֹוָה׃}
}
{
	\eP{
		And Aaron spoke to Moses: “Here, today they brought
		forward their sin offering and their burnt offering in
		front of YHWH. When things like these have happened to me,
		if I had eaten a sin offering today would it be good in
		YHWH’s eyes?”
	}
}
{
}

\Verse{20}
{
	\hP{וַיִּשְׁמַע מֹשֶׁה וַיִּיטַב בְּעֵינָיו׃}
}
{
	\eP{And Moses listened, and it was good in his eyes.}
}
{
}
